\chapter{Introducción}
\label{ch:intro}

% Motivación general de la investigación, hipótesis, objetivos y estructura del informe.
% Máximo 5 páginas.

\section{Motivación} %Listo

La \acrfull{sdp} constituye una de las clases de optimización convexa de mayor relevancia
práctica y teórica.
Sus aplicaciones abarcan desde el diseño de controladores robustos y la síntesis de filtros
en ingeniería de sistemas, hasta la relajación de problemas combinatoriales NP-difíciles
—como corte máximo, coloreo de grafos y programación cuadrática binaria—,
el procesamiento de señales, el aprendizaje automático y la geometría algebraica.
Esta versatilidad ha motivado el desarrollo de una variedad de solvers especializados que
implementan distintos paradigmas algorítmicos: métodos de punto interior primal-dual,
métodos de barrera dual, y métodos de Lagrangiano aumentado orientados a gran escala.

Sin embargo, la pregunta de \textit{qué solver es más adecuado para un problema dado}
permanece abierta.
Los benchmarks tradicionales suelen evaluar el desempeño de los algoritmos mediante
métricas agregadas, como el rendimiento promedio sobre un conjunto predefinido de instancias
de prueba, o mediante análisis de peor caso.
Si bien estos enfoques permiten obtener una visión general del desempeño, también pueden
ocultar patrones relevantes asociados a características específicas de las instancias.
En consecuencia, se pierde capacidad para identificar las fortalezas y debilidades
particulares de cada algoritmo frente a distintas estructuras del problema.

El marco metodológico de \acrfull{isa}, introducido por Smith-Miles y Muñoz, ofrece una
respuesta a esta pregunta.
\acrshort{isa} proyecta las instancias de un problema en un espacio bidimensional definido
por sus características (Features) y, sobre ese espacio, delimita las regiones en que
cada algoritmo exhibe rendimiento superior.
Esto permite identificar de forma visual y cuantitativa las \textit{fortalezas} y
\textit{debilidades} de cada solver, así como detectar regiones del espacio de instancias
que se encuentran subrepresentadas en los benchmarks
disponibles~\cite{smith-milesInstanceSpaceAnalysis2023}.

Una limitación frecuente en los estudios basados en \acrshort{isa} es la escasez de
instancias en ciertas regiones del espacio, lo que dificulta la extracción de conclusiones
robustas sobre el rendimiento de los algoritmos en esas zonas.
Para subsanar esta limitación, Smith-Miles y Bowly~\cite{smith-milesGeneratingNewTest2015}
propusieron el uso de algoritmos genéticos que evolucionan instancias existentes hacia
regiones objetivo del espacio de instancias, generando nuevas instancias con características
deseables de forma automatizada.

La aplicación conjunta de \acrshort{isa} y generación evolutiva de instancias al dominio de
la \acrshort{sdp} representa una contribución original: a la fecha, no existe un estudio
sistemático del espacio de instancias de la \acrshort{sdp} que integre caracterización
por medio de Features, análisis comparativo de solvers y generación dirigida de nuevas instancias.
Este trabajo aborda esa brecha.

\section{Hipótesis}%Listo

Las instancias de \acrshort{sdp} se distribuyen de forma no uniforme en el espacio definido
por sus features, y esta distribución determina en gran medida el
rendimiento relativo de los solvers disponibles. En consecuencia, es posible: (i) caracterizar
las regiones del espacio de instancias en que cada solver presenta ventajas comparativas,
identificando las features que condicionan su desempeño; y (ii) generar mediante
evolución dirigida nuevas instancias con características deseables que amplíen la
cobertura experimental del espacio, enriqueciendo así la capacidad de análisis y de selección de
algoritmos.

\section{Objetivos}

\subsection{Objetivo General}%Listo

Caracterizar el espacio de instancias de la \acrfull{sdp} mediante el análisis sistemático de
features, el estudio comparativo del rendimiento de solvers representativos,
y la generación evolutiva de instancias sintéticas que amplíen la cobertura experimental del espacio de
instancias.

\subsection{Objetivos Específicos}%Listo

\begin{enumerate}
    \item Diseñar un conjunto de features computables directamente desde instancias
    \acrshort{sdp} en formato SDPA sparse, que capturen dimensiones relevantes
    de la geometría del problema tales como la estructura de bloques, la sparsity, la conectividad del grafo
    agregado y las propiedades de chordalidad.

    \item Construir el espacio de instancias de \acrshort{sdp} mediante el framework \acrshort{isa} y
    analizar las fortalezas y debilidades de un conjunto representativo de solvers, identificando las regiones
    del espacio en que cada algoritmo exhibe rendimiento superior y las features que explican
    dicho rendimiento.

    \item Generar instancias sintéticas de \acrshort{sdp} con características estructurales deseables mediante
    un \acrfull{ga} que evoluciona instancias existentes hacia regiones subpobladas del espacio de instancias,
    ampliando la cobertura experimental disponible para el análisis y permitiendo la generación de instancias
    con características deseables para otros fines.
\end{enumerate}

\section{Estructura del informe}%Por hacer

% Descripción breve de la organización del documento.

